% Transcription — fonds Alexandre Grothendieck, shelfmark 71, batch 2 (pages 21-40).
% Established from the facsimile archives/batches/71/batch-02.pdf (a mirror of
% https://grothendieck.umontpellier.fr/71.pdf), per the skill
% transcribe-grothendieck. The batch file carries no cover sheet: PDF page k is
% archivists' page 20+k; the pencilled numbers bottom-left (« 21 » ... « 40 »)
% agree.
%
% Pass: Opus 5.5 (claude-opus-5-5), 2026-09-24 — first pass, unchecked against the pages by a human.
% Revised: Opus 5.5 (claude-opus-5-5), 2026-09-24 — the same day: the mémoire pages are summarised (see below).
%
% Two kinds of page, treated differently.
%   - His own leaves, in the black felt-tip hand of his working notes (and
%     page 29 in pencil): 21, 23, 27, 29, 31. Transcribed in full; 23 and 27
%     carry a single drawing each, described in a note.
%   - A mémoire in another hand (blue ink, paginated in Roman numerals),
%     signed and dated 15 juillet 1976 on page 47, of which this batch holds
%     pages II-XIII: 22 (II), 24 (III), 26 (IV), 28 (V), 30 (VI), 32 (VII),
%     33 (VIII), 35 (IX), 36 (X), 37 (XI), 38 (XII), 39 (XIII). By the user's
%     decision of 2026-09-24, and following folder 100's rule for other
%     people's text (header of transcripts/100/batch-02.fr.tex), these pages
%     are SUMMARISED in one French \note each (the lemmas and definitions
%     they state, named and numbered, one sentence each, no proofs), not
%     recomposed. His marks on them — the black-ink insertions,
%     underlinings, « ! », « style », « vague », « jargon ? », « C'est une
%     salade… », « faux », « ambigu », and the substantive glosses (on
%     Remarque 2, π_0(D_x − D_x∩Γ) → π_0(D'_x − D'_x∩Γ) bijective; on Lemme
%     V, Lemme VII, Lemme XI) — are given in full as \marginal{}, each tied
%     to the words of the mémoire it answers. Pencil remarks on 35-37 (« faux
%     pour n = 3 ! », « il faut prendre s comme dans la dém. du lemme VIII »,
%     « ce n'est pas un “ens. convexe” ! », « a = s ? b = s' ») are of a hand
%     not settled here and are given as \marginal with that caution.
% Absent (no \page): 25, 34 and 40, versos of the mémoire blank but for an
% abandoned start in its author's hand, upside down at the foot. The
% mémoire's author is not named in this edition.
%
% What the batch is. The middle of the mémoire's proof of the polygonal
% Jordan theorem (definitions of J-polygone, bon sommet, visibility; Lemmes
% I-XV; § 2 on convex sets; § 3 « Preuve du théorème 1 » up to the induction
% on card S(Γ) − 2), and his reading of it:
%   21     His notes on the mémoire's Remarque 1: Γ^1_{xy}, Γ^2_{xy}; NB it
%          is the unordered pair, not the couple, that is unique; the
%          operations « rajouter sommets » and division / recollement.
%   29     His pencil sketch towards Lemme XI: Φ := S(Γ) ∩ T° − {s,s',s''},
%          a vertex t of Φ farthest from the line s's'', « Cor Main », and a
%          lemma on a vertex s whose > π sector is exterior, with cases (i),
%          (ii) broken off.
%   31     His lemma: two distinct polygons meet in finitely many arcs;
%          for Γ' = ∂E(Γ), ]s,t[ ⊂ Γ' ∩ ∁Γ ⊂ Ext Γ; Cor.; NB Ext is defined
%          for any polygon.
%
% Notation, fixed once, in his pages and in the summaries. The barred C for
% a convex set is $C$; a small ring over a letter (C°, T°) is \mathring;
% « int », « ext », « Int », « Ext » are \operatorname as written; the empty
% set, drawn φ, is \emptyset; S'(Γ) where the S carries a tick (page 21).
% His insertions in his own leaves are integrated silently; cancellations
% are \struck{}.
%
% The dating is the folder's own; the group « Géométrie et topologie
% combinatoire » (69-89) holds it.
\documentclass[11pt,a4paper]{article}
% Le préambule commun à toutes les transcriptions du fonds Grothendieck.
%
% Il tient en une idée : **l'incertitude se compose, elle ne se résout pas.**
% Une transcription qui lisse un mot douteux a détruit la seule chose pour
% laquelle on la faisait — un lecteur doit pouvoir savoir, à chaque endroit,
% ce qui était sur le feuillet et ce qui a été ajouté. Les six macros qui
% suivent sont donc l'appareil critique, et non de la décoration.
%
% Les mêmes commandes sont reconnues par `scripts/render.mjs`, qui produit la
% vue de lecture du site. Une macro ajoutée ici doit l'être aussi là-bas,
% faute de quoi le rendu échouera bruyamment — ce qui est voulu : un rendu
% silencieusement faux est pire qu'un rendu absent.

\ProvidesPackage{grothendieck}[2026/08/08 Transcriptions du fonds Grothendieck]

\usepackage{amsmath,amssymb,amsthm}
\usepackage{mathrsfs}
\usepackage{stmaryrd}
% Les diagrammes commutatifs sont partout dans ce fonds : c'est la langue
% même de la géométrie algébrique de Grothendieck, et une transcription qui
% les rendrait en prose ne transcrirait rien.
\usepackage{tikz-cd}
% Français par défaut — c'est la langue des feuillets — mais l'anglais est
% chargé aussi : la traduction et le résumé sont des documents anglais, et la
% typographie française (espaces avant les deux-points) y serait une faute.
% Ces documents basculent par \selectlanguage{english} en tête de corps.
\usepackage[english,french]{babel}
\usepackage{fontspec}
\usepackage{geometry}
\geometry{a4paper,margin=25mm,marginparwidth=28mm}
\usepackage{marginnote}
\usepackage{xcolor}
% ulem, not soul. `soul`'s \st and \ul are built on a character-by-character
% scan that XeLaTeX's font machinery breaks: the document fails with an
% undefined control sequence rather than degrading. ulem does the same two jobs
% and is Unicode-engine safe. `normalem` stops it hijacking \emph, which the
% transcriptions use for Grothendieck's own emphasis.
\usepackage[normalem]{ulem}
\usepackage{hyperref}
\hypersetup{colorlinks=true,linkcolor=black,urlcolor=black,pdfborder={0 0 0}}

% --- Métadonnées du lot -----------------------------------------------------
% Reprises telles quelles par le script de rendu, qui les affiche en tête de
% la vue de lecture. Elles fixent ce que le fichier prétend couvrir : sans
% elles, une transcription flotte sans point d'attache dans un fonds de seize
% mille feuillets.

\newcommand{\folder}[1]{\gdef\grothFolder{#1}}
\newcommand{\batch}[1]{\gdef\grothBatch{#1}}
\newcommand{\leaves}[2]{\gdef\grothFirst{#1}\gdef\grothLast{#2}}
\newcommand{\foldertitle}[1]{\gdef\grothFoldertitle{#1}}
\gdef\grothFolder{?}\gdef\grothBatch{?}\gdef\grothFirst{?}\gdef\grothLast{?}\gdef\grothFoldertitle{}

% --- Le résumé --------------------------------------------------------------
%
% Il ouvre la lecture modernisée, sous le titre « Résumé », et il est composé
% en petit. Ce n'est pas un ornement : c'est le seul passage du document qui
% ne soit pas des mathématiques, et un lecteur doit pouvoir le reconnaître
% avant de l'avoir lu — puis le sauter s'il connaît déjà le sujet, ou s'y
% arrêter s'il ne le connaît pas. Le filet qui le suit marque la reprise du
% corps.

\newenvironment{resume}
  {\small\setlength{\parskip}{0.45em}}
  {\par\medskip\hrule\medskip}

% --- Mots-clés ---------------------------------------------------------------
%
% En anglais, en clôture du résumé de la lecture modernisée : le vocabulaire
% moderne sous lequel chercher ce que les feuillets construisent. Le manifeste
% du site (`npm run manifest`) les extrait pour étiqueter la cote entière —
% cette ligne est la source unique des tags, il n'y a pas de fichier à côté.
\newcommand{\keywords}[1]{\par\smallskip\noindent
  {\footnotesize\textsc{Keywords} --- \textit{#1}}\par}

% --- L'appareil critique ----------------------------------------------------

% Le numéro de feuillet, en marge. C'est l'ancre qui permet de régler un
% désaccord sur une formule en regardant *un* feuillet plutôt qu'une chemise
% de six cents. Le site s'en sert aussi pour tourner les pages du fac-similé
% quand on fait défiler la transcription.
\newcommand{\leaf}[1]{%
  \marginnote{\footnotesize\color{gray!70}\textsf{#1}}%
  \label{leaf:#1}\ignorespaces}

% Illisible : on ne devine pas. Un mot inventé qui se lit comme les autres est
% le pire résultat possible.
\newcommand{\ill}{\textcolor{red!60!black}{[\,\ldots\,]}}

% Lecture douteuse : le mot est donné, et signalé comme incertain.
\newcommand{\uncertain}[1]{\uline{#1}}

% Ajout de l'éditeur — une lettre manquante, un mot sous-entendu.
\newcommand{\add}[1]{\textcolor{blue!50!black}{[#1]}}

% Biffé par Grothendieck lui-même. Ce qu'il a rayé fait partie du document :
% une rature dit souvent où la pensée a changé de direction.
\newcommand{\struck}[1]{\sout{#1}}

% Note du transcripteur, sur l'état matériel du feuillet.
\newcommand{\note}[1]{\par\smallskip{\small\color{gray!60!black}\textit{[#1]}}\par\smallskip}

% Note marginale de Grothendieck, distincte du corps du texte.
\newcommand{\marginal}[1]{\par\smallskip\noindent\hspace{1em}%
  {\small\color{gray!60!black}#1}\par\smallskip}

% --- Titre ------------------------------------------------------------------

\AtBeginDocument{%
  \begin{center}
    {\large\bfseries Fonds Alexandre Grothendieck --- cote n\textsuperscript{o}~\grothFolder}\\[2pt]
    {\itshape\grothFoldertitle}\\[4pt]
    {\small feuillets \grothFirst--\grothLast\ (lot \grothBatch)}\\[2pt]
    {\footnotesize\color{gray!60!black}Transcription établie d'après le fac-similé numérisé
     par l'Université de Montpellier.\\
     Les lectures incertaines sont soulignées, les illisibles notées [\,\ldots\,],
     les ajouts entre crochets.}
  \end{center}
  \vspace{1em}}

\endinput


\folder{71}
\batch{2}
\pages{21}{40}
\dating{[à partir de 1976]}
\watermark{Édition de démonstration}
\foldertitle{Théorème de Jordan : notes manuscrites (1976, s.d.).}

\begin{document}

\page{21}
\note{de sa main, à l'encre noire, sur un feuillet séparé. En haut à gauche,
un croquis : un triangle, deux points $x$, $y$ sur sa base, et deux arcs
marqués $1$ et $2$ qui les relient. Les formules reprennent la Remarque 1 du
mémoire, page 20}

\[
\left\{
\begin{array}{l}
\Gamma^1_{xy} \cap \Gamma^2_{xy} = \{x,y\} \\[2pt]
\Gamma^1_{x,y} \cup \Gamma^2_{x,y} = \Gamma \\[2pt]
S'(\Gamma^1_{xy}) \cup S'(\Gamma^2_{xy}) = S(\Gamma) \smallsetminus \{x,y\}
\end{array}
\right.
\]
\note{la dernière ligne se prolonge au bord droit par quelques signes serrés,
peut-être « $\cap\, S(\Gamma)$ », que la présente lecture ne déchiffre pas}

\textbf{NB} Ce n'est pas le \emph{couple} $(\Gamma^1_{xy}, \Gamma^2_{xy})$
qui est déterminé de façon unique, mais l'\emph{ens.}
$\{\Gamma^1_{xy}, \Gamma^2_{xy}\}$.

a) op.\ de \uncertain{rajouter} sommets

b) op.\ de division d'un polygone (par 2 sommets distincts) en deux lignes
polygonales (op.\ inverse : recollement de 2 lignes polygonales ayant
\uncertain{mêmes} extrémités).

\note{plus bas, deux dessins sans lettres : un polygone non convexe en dents
de scie, et un petit triangle dont un côté est tracé en pointillé, avec un
point au-dessus et un point au-dessous de ce côté}

\page{22}
\note{p.\ II du mémoire d'une autre main, résumée et non recomposée. Elle
achève la Remarque 1 ; note qu'on recolle selon $\{x,y\}$ deux arcs
$\Lambda^1_{xy}$, $\Lambda^2_{xy}$ tels que $\Lambda^1_{xy} \cap
\Lambda^2_{xy} = \{x,y\}$ en un polygone $\Lambda$ ; donne la
\emph{Définition} 1 (J-polygone : a) $\mathbb{R}^2 - \Gamma = U \cup V$,
$U$, $V$ connexes, $U$ relativement compact, $V = \operatorname{ext}\Gamma$
non compact ; b) un disque $D_x$ centré sur $\Gamma$ qui ne coupe que les
arêtes contenant $x$ vérifie $D_x - D_x \cap \Gamma = U' \cup V'$, $U'
\subset \operatorname{int}\Gamma$, $V' \subset \operatorname{ext}\Gamma$
connexes) ; la \emph{Remarque} 2 (il suffit qu'un tel disque vérifie b)) ;
$R_\Gamma = \operatorname{int}\Gamma \cup \Gamma$ ; la \emph{Définition} 2
(bon sommet : non colinéaire à ses deux voisins, triangle sans autre sommet
de $\Gamma$ à l'intérieur) ; et, sur une bande collée, la définition de
« $s'$ est $E$-visible de $s$ » ($]s,s'[\, \subset E$) et de la visibilité
de l'intérieur ou de l'extérieur. Ses marques sont à l'encre noire}

\marginal{en face de « on peut obtenir par recollement selon $\{x,y\}$ un
polygone $\Lambda$ » : \uncertain{bien}}

\marginal{en face de « $\overline{V}$ non-compact », qu'il souligne en
tirets : « ! $V$ non rel.\ \uncertain{compact} »}

\marginal{après « $V' \subset \operatorname{ext}\Gamma$ connexes » :
« non vides ! », et sous la ligne : « (i.e.\ les 2 comp.\ conn.\ de $D_x -
D_x \cap \Gamma$ \struck{\uncertain{appart.}}, l'une dans
$\operatorname{int}\Gamma$, l'autre dans $\operatorname{Ext}\Gamma$) »}

\marginal{en face de « Remarque 2 » : « style »}

\marginal{à la fin de la Remarque 2, « la vérifie » : « (car si $D_x
\subset D'_x$, $\pi_0(D_x - D_x \cap \Gamma) \to \pi_0(D'_x - D'_x \cap
\Gamma)$ est bijectif) »}

\marginal{il souligne « bon sommet » (Définition 2)}

\marginal{en tête de la bande collée : « vague », et « Déf 3 : » ; il
insère « et $s, s' \in \mathbb{R}^2$, » après « Si $E$ est une partie de
$\mathbb{R}^2$ » et biffe la même clause en fin de phrase}

\page{23}
\note{sans texte : un seul dessin, à l'encre noire, en haut à gauche — un
polygone étoilé non convexe, doublé à l'intérieur d'un second contour
parallèle, plus fin}

\page{24}
\note{p.\ III du mémoire. \emph{Lemme} I : si $\Gamma$ est un J-polygone et
$\Gamma'_{xy}$ l'un des deux arcs de la Remarque 1 reliant $x \neq y$, il
existe un système fondamental $\beta$ de voisinages ouverts $V$ de
$\Gamma'_{xy}$ tels que $V \cap \operatorname{int}\Gamma$ et $V \cap
\operatorname{ext}\Gamma$ soient connexes. Suit la démonstration, par un
recouvrement fini de l'arc par des disques vérifiant b) de la définition 1}

\marginal{il insère « $x \neq y$ », « $\beta$ » et « ouverts » dans
l'énoncé, et souligne « il existe un système » ; « ! » en marge}

\marginal{en face de « grâce à l'axiome du choix », qu'il souligne : « ! » ;
il souligne « de $\Gamma$ » dans « si $U$ est un voisinage de $\Gamma$ » et
écrit en marge : « $\Gamma_{xy}$ »}

\marginal{il encadre « i.e.\ $\Gamma_{xy} \cap D_{z_i} = I_i$ »}

\marginal{en face de « Le même argument prouve (quitte à changer
$\operatorname{int}\Gamma$ par $\operatorname{ext}\Gamma$) », qu'il
souligne : « ! »}

\marginal{en face de « On utilise le fait, sans le dire, que deux points
$x, y$ d'un ouvert connexe de $\mathbb{R}^2$ peuvent être reliés par un arc
de polygone », où il souligne « sans le dire » : « ? »}

\page{26}
\note{p.\ IV du mémoire : « 2 — Rappel sur les sous-ensembles convexes de
$\mathbb{R}^2$ ». \emph{Lemme} II : si $A \subset \mathbb{R}^2$ contient
trois points affinement indépendants, il existe un unique polygone $\Gamma$
sans trois sommets consécutifs colinéaires, d'ensemble sous-jacent $E(A) -
\mathring{E}(A)$ ($E$ = enveloppe convexe), et $S(\Gamma) \subset A$.
\emph{Convention} : « ensemble convexe » désigne l'enveloppe convexe d'un
ensemble fini de points ; si $\mathring{C} \neq \emptyset$, $\partial C = C -
\mathring{C}$ reçoit la structure de polygone du lemme II. \emph{Lemme} III :
la droite de deux sommets consécutifs de $\partial C$ laisse $\mathring{C}$
d'un seul côté. \emph{Lemme} IV : si $x \notin C$, il existe $y \in C -
S(\partial C)$ qui est $(\mathbb{R}^2 - C)$-visible de $x$. Démonstrations
esquissées}

\marginal{en face de « Lemme II » : « ! » ; il insère « finie » après « une
partie » et « de $\Gamma$ » après « trois sommets consécutifs »}

\marginal{en face de « Démonstration (Esquisse) » : « jargon ? » ; il entoure
le $\bigcap$ de « $E(A) = \bigcap$ demi-espaces fermés », souligne
« uniquement déterminées par » et insère « ens.\ des » dans « $S(\Gamma)$ =
points extrémaux de $E(A)$ »}

\marginal{en face de « Convention » : « !!! » et
« \uncertain{partie polyédr.\ conv.} » ; il souligne « pour nous » et
« ensemble convexe »}

\marginal{à droite de « un ensemble convexe $C$ est compact » :
« intervention \uncertain{en} lemme II ! C'est une salade… »}

\marginal{en face de la démonstration du lemme III : « ? »}

\marginal{en face de « Lemme IV » : « vis. » ; en face de sa dernière ligne,
« la distance de $x$ à $[z_1, z_2]$ » : « ! »}

\page{27}
\note{sans texte : un croquis à l'encre noire — un segment d'extrémités $s$
et $s'$, un point $y$ sur ce segment, et une droite en pointillé qui monte de
$y$ vers un point marqué $x$ ; il paraît illustrer la démonstration du lemme
IV, page 26}

\page{28}
\note{p.\ V du mémoire. \emph{Lemme} V : si $C$ est convexe, $y \in\,
]s,s'[$ avec $s, s' \in C$, et $y$ est $(\mathbb{R}^2 - C)$-visible d'un point
$x \notin C$ non colinéaire à $s, s'$, alors tout point colinéaire à $s, s'$
est $(\mathbb{R}^2 - C)$-visible de $x$ (démonstration par le triangle $s,
s', w$). \emph{Définition} 3 : un polygone est convexe s'il est le bord
$C - \mathring{C}$ d'un ensemble convexe. \emph{Lemme} VI : un polygone convexe
est un J-polygone, et tout sommet non colinéaire à ses deux voisins est un
bon sommet}

\marginal{au-dessus de « Lemme V » : « vis ? » ; il souligne « ensemble
convexe »}

\marginal{entre l'énoncé du lemme V et sa démonstration : « \textbf{NB}
$s, s'$ appartiennent : \ill{} \ill{} \uncertain{arête} de $C$, et si $D$
est la droite qu'ils déterminent, $x$ et $C$ sont de part et d'autre de
cette droite… »}

\marginal{au-dessus de « Définition 3. Un polygone $\Gamma$ est convexe » :
« (\uncertain{pas} \ill{} de \uncertain{jargon} !) » ; il insère « dit »
avant « convexe » et souligne en tirets « ensemble »}

\page{29}
\note{verso du feuillet de la page 28, de sa main, au crayon ; ce qui suit
jusqu'au « Cor Main » est écrit en travers du feuillet, le lemme du bas dans
le sens de la page. Deux dessins : au crayon, un polygone en dents de scie
inscrit dans un grand triangle ; à l'encre noire, un polygone en pointe de
flèche dont deux côtés sont complétés en pointillé}

\struck{Or [les comp.\ connexes de $\Gamma \cap [s', s'']$ sont des
\ill{} des points, et les extrémités des $I_\alpha$ et] \ill{} côtés
\ill{} et que tout $t \in$}
\note{passage barré de grands traits obliques ; la lecture en est partielle}

alors $\Phi := S(\Gamma) \cap \mathring{T} - \{s, s', s''\}$
\struck{\ill{}} $\neq \emptyset$, et $t \in \Phi$ réalise le maximum de la
distance à la droite définie par $s', s''$, alors $]s,t[\, \subset
\operatorname{Int}\Gamma$.

\emph{Cor Main.} Ou bien $]s', s''[\, \subset \operatorname{Int}\Gamma$, ou
bien $\exists\, t \in S(\Gamma)$, $t \neq s$, tel que $]s,t[\, \subset
\operatorname{Int}\Gamma$.

\emph{Lemme}. Soit $\Gamma$ un polygone à $n$ côtés ($n \geq 3$)
\struck{non convexe}. Alors $\exists\, s \in S(\Gamma)$ « non plat » tel que
$\forall D_s$ de centre $s$, \struck{l'un des} satisfaisant la cond.\ b) de
Déf 1, le secteur angulaire de $D$ découpé par $\Gamma$ (i.e.\ par $[s,s']
\cup [s,s'']$, où $s', s''$ sont les 2 sommets adjacents) qui est d'angle $>
\pi$ soit extérieur à $\Gamma$. Si $\Gamma$ est un J-polygone, et $n \geq 4$,
alors

(i) si $\Gamma \cap \mathring{T} = [s,s'] \cup [s,s'']$, $]s',s''[\,
\subset \operatorname{Int}\Gamma$ \struck{\ill{}}

(ii) si $\Gamma \cap T \neq [s,s'] \cup [s,s'']$ \struck{et $\Gamma \cap
\mathring{T} = \emptyset$, \ill{} $]s', s''[$, \ill{}}
\note{la page s'arrête là. Le $\mathring{T}$ de (i) est récrit sur un mot
noirci ; $T$ est le triangle de sommets $s, s', s''$, comme au lemme XI de
la page 35}

\page{30}
\note{p.\ VI du mémoire : démonstration du lemme VI (par un triangle $T$
contenant $C$ et la connexité de $(\partial T \cup
\operatorname{int}\partial T) - C$). Puis, pour un polygone $\Gamma$ et son
enveloppe convexe $E(\Gamma)$ : a) $E(\Gamma) = E(S(\Gamma))$ ; b)
$S(\partial(\Gamma)) \subset S(\Gamma)$. « On veut étudier le lien entre
$\Gamma$ et son enveloppe convexe. »}

\marginal{en face de « Il est clair que $E(\Gamma)$ vérifie », où il
souligne « clair » : « \uncertain{réf.\ int.} »}

\marginal{sous « b) $S(\partial(\Gamma))$ » : « $S(\partial(E(\Gamma)))$ »}

\page{31}
\note{de sa main, à l'encre noire, sur un feuillet séparé. À gauche, un
dessin : un polygone non convexe à six côtés, en éclair}

\emph{Lemme}. Soient $\Gamma, \Gamma'$ deux polygones distincts,
\struck{tels que $\Gamma \subset$} alors $\Gamma \not\subset \Gamma'$,
$\Gamma' \not\subset \Gamma$, et $\Gamma \cap \Gamma'$ est une réunion finie
d'arcs disjoints de $\Gamma$ (et de $\Gamma'$), éventuellement réduits à un
point. Si $\Gamma' = \partial E(\Gamma)$, où $\Gamma$ est un polygone
\uncertain{non convexe} (puisque $\Gamma' \neq \Gamma$), alors
\struck{les} l'\uncertain{ens.\ $\Phi$ des} extrémités des arcs polygonaux
composantes connexes de $\Gamma \cap \Gamma'$ \struck{sont des sommets} est
$\subset S(\Gamma)$\struck{\ill{}}. Alors pour toute \struck{telle} $s \in
\Phi$, $\exists\, t \in \Phi$, $t \neq s$, telle que $]s,t[\, \subset
\Gamma'$ \struck{\ill{}} $\cap\, \complement\Gamma$ et alors $]s,t[\,
\subset \operatorname{Ext}\Gamma$.
\marginal{\textbf{NB} $t$ est unique, sauf si $s$ point isolé de $\Gamma
\cap \Gamma'$, auquel cas il existe exactement deux tels $t$.}

\emph{Cor}. Si $\Gamma$ polygone non convexe, $\exists\, s, s'$ sommets
distincts de $\Gamma$, tels que $s$ \struck{\ill{}} $]s,s'[\, \subset
\operatorname{Ext}\Gamma \cap \partial E(\Gamma)$.

[\textbf{NB} $\operatorname{Ext}(\Gamma)$ est défini pour \uncertain{tt}
polygone, pas \uncertain{néc.}\ un J-polygone…]

\page{32}
\note{p.\ VII du mémoire. \emph{Lemme} VII : si $\Gamma$ est un J-polygone,
$\operatorname{ext}\partial E(\Gamma) \subset \operatorname{ext}\Gamma$.
\emph{Lemme} VIII : si $\Gamma$ est un J-polygone non convexe, il existe deux
sommets $s, s'$ tels que $s'$ soit visible de $s$ de l'extérieur de $\Gamma$ ;
la démonstration prend une arête $\alpha = [s,s']$ de $\partial E(\Gamma)$
non contenue dans $\Gamma$, et se poursuit page 33}

\marginal{en face de « Lemme VIII » : « \uncertain{inutile} » ; il met entre
parenthèses le « J- » de « un J-polygone non-convexe »}

\marginal{en face de la démonstration du lemme VIII, une note verticale
partiellement lue : « \ill{} $E(\Gamma)^c$ \ill{} $E(\partial E(\Gamma)) =
E(\Gamma)$ \ill{} »}

\page{33}
\note{p.\ VIII du mémoire : fin de la démonstration du lemme VIII
($]s,s'[\, \subset \operatorname{ext}\Gamma$). Secteurs angulaires intérieur
et extérieur en un sommet d'un J-polygone. \emph{Lemme} IX (« l'analogue
suivant de la positivité de la courbure des courbes convexes ») : en un
sommet d'un polygone convexe, l'angle du secteur intérieur est $\leq \pi$.
\emph{Lemme} X : un J-polygone a au moins trois sommets où l'angle du
secteur intérieur est $< \pi$ ; la démonstration part d'un sommet de
$\partial E(\Gamma)$ et se poursuit page 35}

\marginal{à droite de « Mais $D_t \cap \operatorname{ext}\partial E(\Gamma)
\neq \emptyset$ », encadré : « donc $D_t \cap \operatorname{ext}\Gamma \neq
\emptyset$ (cf lemme VII) »}

\note{le « J- » de l'énoncé du lemme X est barré d'un trait de crayon, d'une
main non établie}

\page{35}
\note{p.\ IX du mémoire : fin de la démonstration du lemme X (l'angle
extérieur en un sommet de $\partial E(\Gamma)$ est $> \pi$ et reste dans
$\operatorname{ext}\Gamma$). \emph{Lemme} XI : si $\Gamma$ est un J-polygone,
il existe deux sommets $s, s'$ tels que $s'$ soit visible de $s$ de
l'intérieur de $\Gamma$. La démonstration prend $s$ comme au lemme VIII, ses
voisins $s', s''$ et le triangle $T$ de sommets $s, s', s''$, et traite
d'abord le cas $(T - ([s,s'] \cup [s,s''])) \cap \Gamma = \emptyset$. Les
marques au crayon sont d'une main que la présente lecture ne distingue pas
avec certitude de celle de l'encre noire}

\marginal{en face de « Lemme XI », au crayon : « faux pour $n = 3$ ! » ;
à l'encre noire, il insère « à $n \geq 3$ côtés, » après « J-polygone » et
écrit sous l'énoncé : « OPS $\Gamma$ non convexe (sinon c'est clair grâce à
…) »}

\marginal{en face de « Soit $s$ un sommet satisfaisant la propriété
énoncée au lemme VIII » : « Lemme X » ; en face de « Si $(T - ([s,s'] \cup
[s,s''])) \cap \Gamma = \emptyset$ » : « ambigu »}

\marginal{en face de « donc $]s',s''[\, \subset \operatorname{int}\Gamma$ » :
« faux » ; au crayon, à la suite : « il faut prendre $s$ comme dans la
\emph{dém.}\ du lemme VIII (pas son énoncé !) »}

\marginal{en face de « Si $(T - ([s,s'] \cup [s,s''])) \cap \Gamma \subset\,
]s',s''[$ », où il barre le $\subset$ et souligne la ligne : « \struck{\ill{}}
! » et, entouré, « non vide ? »}

\marginal{au crayon, sous « $\mathring{T} \subset \operatorname{int}\Gamma$ »
entouré : « dire avant ! (\ill{} au lemme VIII) »}

\page{36}
\note{p.\ X du mémoire : fin de la démonstration du lemme XI, dans le cas
$\mathring{T} \cap \Gamma \neq \emptyset$, par l'enveloppe convexe $C$ de
$(\mathring{T} \cup\, ]s', s''[) \cap \Gamma$ et les lemmes IV et V.
\emph{Corollaire à la démonstration du lemme VIII} : tout polygone a deux
sommets distincts $s, s'$ avec $s'$ $(\mathbb{R}^2 - \Gamma)$-visible de $s$
et $]s', s[\, \subset \partial E(\Gamma)$}

\marginal{en face de « Reste le cas $\mathring{T} \cap \Gamma \neq
\emptyset$ » : « !! » ; au crayon, au-dessus de « l'enveloppe convexe $C$ de
$(\mathring{T} \cup\, ]s',s''[) \cap \Gamma$ » : « ce n'est pas un
“ens.\ convexe” ! » ; à l'encre noire, des crochets tracés par-dessus font
de $]s',s''[$ un $[s',s'']$}

\marginal{au crayon, dans la marge gauche, en regard de « Soit $\alpha \in
A(\partial C)$ » et de la suite, une longue note verticale : « \ill{}
\uncertain{compliqué} \ill{} $\in \Gamma$ \ill{} $(s', s'')$ »}

\marginal{après « Corollaire à la démonstration du lemme VIII », qu'une
main souligne au crayon : « !! » ; au crayon, « distincts » inséré avant
« $s, s' \in S(\Gamma)$ »}

\page{37}
\note{p.\ XI du mémoire : « 3 — Preuve du théorème 1 ». On prend $s, s'$ comme
au corollaire, les deux arcs $\Gamma^1_{ss'}$, $\Gamma^2_{ss'}$ de $\Gamma$
reliant $s$ à $s'$, et les polygones $\Lambda^1$, $\Lambda^2$ obtenus en
recollant $[s,s']$ à chacun. \emph{Lemme} XII : si $\Lambda^1$ est un
J-polygone, $\Gamma^2_{ab} - \{a,b\}$ est contenu dans
$\operatorname{int}\Lambda^1$ ou dans $\operatorname{ext}\Lambda^1$ (et de
même en échangeant les indices). On suppose ensuite $\Lambda^1$, $\Lambda^2$
J-polygones. \emph{Lemme} XIII : pour $x \in \partial E(\Gamma)$ et un disque
$D_x$ convenable, $\operatorname{ext}\partial E(\Gamma) \cap D_x \subset
\operatorname{ext}\Lambda^i$, $i = 1, 2$}

\marginal{au crayon, en face de « Lemme XII », qui écrit $\Gamma^2_{ab}$ :
« $a = s$ ? $b = s'$ »}

\marginal{au crayon, en face de l'énoncé du lemme XIII, où le $D_x$ de la
formule est barré, encadré : « $\operatorname{Ext}(\partial E(\Gamma)) =
\mathbb{R}^2 - E(\Gamma) \subset \operatorname{Ext}(\Lambda^i)$ »}

\page{38}
\note{p.\ XII du mémoire, sans annotation. $\Lambda^1 \cap \Lambda^2 =
[s,s']$. \emph{Lemme} XIV : pour un disque $D_s$ convenable,
$\operatorname{int}\Lambda^1 \supset \Lambda^2 \cap D_s - [s,s'] \cap D_s$ ou
$\operatorname{int}\Lambda^2 \supset \Lambda^1 \cap D_s - [s,s'] \cap D_s$.
D'où $\operatorname{int}\Lambda^1 \supset \Lambda^2 - \{s,s'\}$ ou
$\operatorname{int}\Lambda^2 \supset \Lambda^1 - \{s,s'\}$ ; on suppose le
premier cas. \emph{Lemme} XV : $\operatorname{int}\Lambda^2 \subset
\operatorname{int}\Lambda^1$}

\page{39}
\note{p.\ XIII du mémoire, sans annotation. Démonstration du lemme XV ; puis
la récurrence sur $\operatorname{card} S(\Gamma) - 2$ (le cas du triangle
étant trivial), avec $\operatorname{card} S(\Lambda^i) <
\operatorname{card} S(\Gamma)$, et les ensembles $V =
\operatorname{ext}\Lambda^1 \cup \operatorname{int}\Lambda^2 \cup\, ]s,s'[$
et $U = \operatorname{int}\Lambda^1 - (\operatorname{int}\Lambda^2 \cup
(\Gamma^2_{ss'} \cap \operatorname{int}\Lambda^2))$, $U$ relativement
compact et $V$ non, que reprend le lemme XVI de la page 41}

\end{document}
